\chapter{Plan de trabajo}	
\label{cap:plantrabajo}

	En este capitulo se especifica el plan de trabajo empelado para el proyecto


\section{Metodología de Trabajo}

\subsection{Marco Metodológico: Scrum Adaptado}
El proyecto se desarrolla bajo un enfoque ágil basado en Scrum, adaptado a las necesidades del equipo y las características del proyecto académico.

\subsubsection{Sprints}
\begin{itemize}
    \item \textbf{Duración:} 2 semanas por sprint
    \item \textbf{Total de sprints:} 6 sprints
    \item \textbf{Duración total:} 12 semanas
\end{itemize}

\subsubsection{Ceremonias Scrum}
\begin{longtable}{|p{3cm}|p{3cm}|p{7cm}|}
\hline
\textbf{Ceremonia} & \textbf{Frecuencia} & \textbf{Descripción} \\
\hline
\endfirsthead

\multicolumn{3}{c}%
{\tablename\ \thetable\ -- \textit{Continuación de la página anterior}} \\
\hline
\textbf{Ceremonia} & \textbf{Frecuencia} & \textbf{Descripción} \\
\hline
\endhead

\hline \multicolumn{3}{r}{\textit{Continúa en la siguiente página}} \\
\endfoot

\hline
\endlastfoot

Sprint Planning & Inicio de cada sprint & Planificación de tareas, definición de objetivos y asignación de historias de usuario \\
\hline
Daily Standup & Lunes, Miércoles, Viernes & Sincronización rápida (15 min) sobre progreso, obstáculos y planes \\
\hline
Sprint Review & Final de cada sprint & Demostración de funcionalidades completadas a stakeholders \\
\hline
Sprint Retrospective & Final de cada sprint & Análisis de mejoras y lecciones aprendidas \\
\hline
\end{longtable}

\subsection{Control de Versiones}
\begin{itemize}
    \item \textbf{Sistema:} Git con GitHub
    \item \textbf{Estrategia de branching:} Git Flow
    \item \textbf{Ramas principales:}
    \begin{itemize}
        \item \texttt{main}: Código en producción
        \item \texttt{develop}: Código en desarrollo
        \item \texttt{feature/*}: Nuevas funcionalidades
        \item \texttt{hotfix/*}: Correcciones urgentes
    \end{itemize}
    \item \textbf{Code Review:} Pull Requests obligatorios con al menos 1 aprobación
\end{itemize}

\subsection{Estándares de Desarrollo}
\begin{itemize}
    \item Código en inglés (variables, funciones, comentarios)
    \item Documentación en español
    \item TypeScript strict mode habilitado
    \item Convenciones de nombres: camelCase para variables, PascalCase para componentes
    \item Commits siguiendo Conventional Commits
\end{itemize}

\section{Organización del Equipo}

\subsection{Roles y Responsabilidades}

\begin{longtable}{|p{3cm}|p{4cm}|p{6cm}|}
\hline
\textbf{Rol} & \textbf{Miembro(s)} & \textbf{Responsabilidades} \\
\hline
\endfirsthead

\multicolumn{3}{c}%
{\tablename\ \thetable\ -- \textit{Continuación de la página anterior}} \\
\hline
\textbf{Rol} & \textbf{Miembro(s)} & \textbf{Responsabilidades} \\
\hline
\endhead

\hline \multicolumn{3}{r}{\textit{Continúa en la siguiente página}} \\
\endfoot

\hline
\endlastfoot

Product Owner & [Aaron Ugalde Téllez] & Definir requisitos, priorizar backlog, validar entregables, comunicación con stakeholders \\
\hline
Scrum Master & [Beltran Saucedo Axel Alejandro] & Facilitar ceremonias, resolver impedimentos, asegurar seguimiento de Scrum \\
\hline
Backend Lead & [Castillo Aldaba Jared] & Diseño de API, arquitectura del servidor, base de datos, seguridad \\
\hline
Frontend Lead & [Vazquez Villagran Jorge] & Arquitectura del cliente, diseño de componentes, UX/UI, integración con API \\
\hline
Desarrollador Full-Stack & [Ugalde Tellez Aaron] & Implementación de features, testing, documentación técnica \\
\hline
QA/Tester & [Beltran Saucedo Axel Alejandro] & Pruebas funcionales, pruebas de integración, documentación de bugs \\
\hline
DevOps & [Castillo Aldaba Jared] & Configuración de entornos, CI/CD, deployment, monitoreo \\
\hline
\end{longtable}

\subsection{Comunicación del Equipo}
\begin{itemize}
    \item \textbf{Herramienta principal:} Discord/Slack
    \item \textbf{Reuniones presenciales:} 2 veces por semana
    \item \textbf{Documentación:} Notion/Confluence
    \item \textbf{Gestión de tareas:} Jira/Trello
    \item \textbf{Repositorio:} GitHub
\end{itemize}

% ============================================
% CRONOGRAMA DE ACTIVIDADES
% ============================================
\section{Cronograma de Actividades}

\subsection{Sprint 0: Planificación e Infraestructura (Semanas 1-2)}

\begin{longtable}{|p{1.5cm}|p{5cm}|p{6.5cm}|}
\hline
\textbf{Día} & \textbf{Actividad} & \textbf{Entregables} \\
\hline
\endfirsthead

\multicolumn{3}{c}%
{\tablename\ \thetable\ -- \textit{Continuación de la página anterior}} \\
\hline
\textbf{Día} & \textbf{Actividad} & \textbf{Entregables} \\
\hline
\endhead

\hline \multicolumn{3}{r}{\textit{Continúa en la siguiente página}} \\
\endfoot

\hline
\endlastfoot

1-2 & Reunión de kick-off & Documento de objetivos y alcance \\
\hline
3-4 & Análisis de requisitos & Documento de requisitos funcionales y no funcionales \\
\hline
5-7 & Diseño de arquitectura & Diagrama de arquitectura, modelo de datos \\
\hline
8-9 & Configuración de repositorio & Repo configurado con estructura de monorepo \\
\hline
10-12 & Setup de entornos & Entorno de desarrollo local configurado \\
\hline
13-14 & Configuración de base de datos & Schema inicial de PostgreSQL \\
\hline
\end{longtable}

\subsection{Sprint 1: Autenticación y Propietarios (Semanas 3-4)}

\textbf{Objetivo:} Implementar el sistema de autenticación y gestión básica de propietarios.

\begin{itemize}
    \item Sistema de registro de propietarios
    \item Login con JWT
    \item Middleware de autenticación
    \item Panel de propietario básico
    \item Gestión de perfil
\end{itemize}

\textbf{Historias de Usuario:}
\begin{itemize}
    \item Como propietario, quiero registrarme en el sistema
    \item Como propietario, quiero iniciar sesión de forma segura
    \item Como propietario, quiero ver y editar mi perfil
\end{itemize}

\subsection{Sprint 2: Gestión de Mascotas (Semanas 5-6)}

\textbf{Objetivo:} Desarrollar el módulo completo de gestión de mascotas.

\begin{itemize}
    \item CRUD de mascotas
    \item Formulario de registro con catálogos (razas, colores, especies)
    \item Listado de mascotas del propietario
    \item Vista detallada de mascota
    \item Validaciones frontend y backend
\end{itemize}

\textbf{Historias de Usuario:}
\begin{itemize}
    \item Como propietario, quiero registrar una mascota
    \item Como propietario, quiero ver todas mis mascotas
    \item Como propietario, quiero editar los datos de mi mascota
\end{itemize}

\subsection{Sprint 3: Historial Médico (Semanas 7-8)}

\textbf{Objetivo:} Implementar el sistema de historial médico completo.

\begin{itemize}
    \item Módulo de vacunas
    \item Módulo de enfermedades
    \item Módulo de alergias
    \item Módulo de desparasitaciones
    \item Subida y gestión de documentos médicos
    \item Vista consolidada de historial
\end{itemize}

\textbf{Historias de Usuario:}
\begin{itemize}
    \item Como propietario, quiero registrar las vacunas de mi mascota
    \item Como propietario, quiero llevar un registro de enfermedades
    \item Como propietario, quiero subir documentos médicos
\end{itemize}

\subsection{Sprint 4: Sistema de Reservaciones (Semanas 9-10)}

\textbf{Objetivo:} Desarrollar el sistema completo de reservaciones.

\begin{itemize}
    \item Catálogo de habitaciones
    \item Creación de reservaciones
    \item Sistema de disponibilidad
    \item Asignación de mascotas a habitaciones
    \item Estados de reservación (pendiente, confirmada, completada)
    \item Cancelación de reservaciones
\end{itemize}

\textbf{Historias de Usuario:}
\begin{itemize}
    \item Como propietario, quiero hacer una reservación para mi mascota
    \item Como propietario, quiero ver mis reservaciones activas
    \item Como admin, quiero confirmar reservaciones
\end{itemize}

\subsection{Sprint 5: Panel Administrativo y Servicios (Semanas 11-12)}

\textbf{Objetivo:} Implementar funcionalidades administrativas y servicios adicionales.

\begin{itemize}
    \item Dashboard administrativo
    \item Gestión de empleados
    \item Catálogo de servicios
    \item Sistema de citas para servicios
    \item Módulo de pagos
    \item Reportes básicos
\end{itemize}

\textbf{Historias de Usuario:}
\begin{itemize}
    \item Como admin, quiero ver todas las reservaciones del hotel
    \item Como admin, quiero gestionar servicios adicionales
    \item Como admin, quiero registrar pagos
\end{itemize}

\subsection{Sprint 6: Testing, Optimización y Deployment (Semanas 13-14)}

\textbf{Objetivo:} Asegurar la calidad del sistema y prepararlo para producción.

\begin{itemize}
    \item Testing exhaustivo de todas las funcionalidades
    \item Corrección de bugs identificados
    \item Optimización de performance (queries, carga de imágenes)
    \item Documentación técnica completa
    \item Documentación de usuario
    \item Deployment a producción
    \item Capacitación al cliente
\end{itemize}

\subsection{Diagrama de Gantt}

\begin{center}
\small
\begin{longtable}{|p{3.5cm}|c|c|c|c|c|c|c|c|c|c|c|c|c|c|}
\hline
\textbf{Actividad} & \multicolumn{14}{c|}{\textbf{Semanas}} \\
\cline{2-15}
& 1 & 2 & 3 & 4 & 5 & 6 & 7 & 8 & 9 & 10 & 11 & 12 & 13 & 14 \\
\hline
\endfirsthead

\multicolumn{15}{c}%
{\tablename\ \thetable\ -- \textit{Continuación de la página anterior}} \\
\hline
\textbf{Actividad} & \multicolumn{14}{c|}{\textbf{Semanas}} \\
\cline{2-15}
& 1 & 2 & 3 & 4 & 5 & 6 & 7 & 8 & 9 & 10 & 11 & 12 & 13 & 14 \\
\hline
\endhead

\hline \multicolumn{15}{r}{\textit{Continúa en la siguiente página}} \\
\endfoot

\hline
\endlastfoot

Sprint 0 & X & X & & & & & & & & & & & & \\
\hline
Sprint 1 & & & X & X & & & & & & & & & & \\
\hline
Sprint 2 & & & & & X & X & & & & & & & & \\
\hline
Sprint 3 & & & & & & & X & X & & & & & & \\
\hline
Sprint 4 & & & & & & & & & X & X & & & & \\
\hline
Sprint 5 & & & & & & & & & & & X & X & & \\
\hline
Sprint 6 & & & & & & & & & & & & & X & X \\
\hline
\end{longtable}
\end{center}
